\documentclass[11pt]{article}
\usepackage[T1]{fontenc}
\usepackage{lmodern}
\usepackage{amsmath,amssymb,amsthm}
\usepackage[margin=1in]{geometry}
\usepackage{microtype}
\usepackage[colorlinks=true,linkcolor=blue,citecolor=blue,urlcolor=blue]{hyperref}
\newtheorem{theorem}{Theorem}[section]
\newtheorem{lemma}[theorem]{Lemma}
\newtheorem{proposition}[theorem]{Proposition}
\theoremstyle{remark}
\newtheorem{remark}[theorem]{Remark}
\newcommand{\R}{\mathbb R}
\newcommand{\Pp}{\mathbb P}
\newcommand{\E}{\mathbb E}
\newcommand{\dist}{\operatorname{dist}}
\newcommand{\diam}{\operatorname{diam}}
\newcommand{\aff}{\operatorname{aff}}
\newcommand{\relint}{\operatorname{relint}}
\newcommand{\vol}{\operatorname{vol}}
\title{Connectivity when isolated vertices disappear\newline in a convex polytope}
\author{Kaiwen Jin\\\small\href{mailto:yui.ui.siki@gmail.com}{yui.ui.siki@gmail.com}}
\date{October 14, 2026}
\begin{document}
\maketitle
\begin{abstract}
Let $X_1,\ldots,X_n$ be independent uniform points in a fixed compact convex
polytope $A\subset\R^d$ with nonempty interior, where $d\geq3$. Let $L_n$ be
the largest nearest-neighbour distance and let $M_n$ be the connectivity
radius of their Euclidean random geometric graph. We prove that
$\Pp(L_n=M_n)\to1$. The argument retains the contribution of every face at a
deterministic radius slightly below the leading connectivity scale. A linear
gain in cone-ball volume excludes small components, and a positive collar
surplus excludes medium separated sets. A fixed spatial net then excludes two
large components simultaneously at every larger radius. An isolated vertex
persists at the chosen lower radius, so these estimates apply at the random
radius $L_n$.
\end{abstract}

\section{Introduction}
For a finite set $\mathcal X\subset\R^d$ and $r\geq0$, let $G(\mathcal X,r)$
join two distinct points when their Euclidean distance is at most $r$.
Suppose that $A\subset\R^d$ is a compact convex polytope with nonempty interior
and that $X_1,X_2,\ldots$ are independent and uniform in $A$. Write
$\mathcal X_n=\{X_1,\ldots,X_n\}$ and, for $n\geq2$, put
\begin{equation}\label{eq:thresholds}
 L_n=\max_{1\leq i\leq n}\min_{j\ne i}\|X_i-X_j\|,
 \qquad
 M_n=\inf\{r\geq0:G(\mathcal X_n,r)\text{ is connected}\}.
\end{equation}
The sample points are distinct almost surely. The first quantity is the
radius at which the last isolated vertex disappears. The second is the
radius at which all components merge. Always $L_n\leq M_n$.

\begin{theorem}\label{thm:main}
For every fixed integer $d\geq3$ and every fixed compact convex polytope
$A\subset\R^d$ with nonempty interior, the uniform sample satisfies
\[
 \lim_{n\to\infty}\Pp(L_n=M_n)=1.
\]
\end{theorem}

Penrose, Yang and Higgs \cite{PYH} obtained a common almost-sure leading
scale for the largest nearest-neighbour link and the connectivity threshold
in a polytope. That scale depends on the dimensions and solid angles of its
faces. Their Remark~2.9 distinguishes this result from exact equality of the
two thresholds with high probability for a general polytope. Penrose and
Yang \cite[Theorem~1.1 and equation~(1.5)]{PY} prove such equality for smooth
domains and, in dimension two, convex polygons. Theorem~\ref{thm:main}
addresses uniform samples in arbitrary fixed convex polytopes of dimension
at least three. It makes no claim about nonuniform densities, higher
connectivity, or uniformity over varying polytopes.

We choose a deterministic radius $r_n$ just below the leading scale and
prove four estimates. At $r_n$ there are no small non-singleton components,
but there is still an isolated vertex. A face-by-face empty-shell estimate
rules out medium separated subsets. A separator and a fixed net rule out two
large components at every radius $s\geq r_n$. These estimates imply that a
disconnected graph at any such radius has an isolated vertex. Applying this
statement at $s=L_n$ proves the theorem without a limit law for $L_n$.

\section{Face geometry and cone volume}
Throughout, $B_r(x)$ is the closed Euclidean ball. We use $|E|$ for
$d$-dimensional Lebesgue measure, $\|x\|$ for Euclidean norm, and natural
logarithms. Set $f_0=|A|^{-1}$, and write
$U_r=(U+B_r(0))\cap A$ for $U\subset A$, where $+$ denotes Minkowski addition.
Constants may depend on the fixed polytope, the dimension, and auxiliary
parameters explicitly fixed in a lemma. They do not depend on $n$ or on a
radius in a stated uniform range. We write $a_n\asymp b_n$ when their ratio is
bounded above and below by positive constants for all sufficiently large $n$.

Let $F$ range over all nonempty faces, including $A$ itself. Put
\[
 k_F=\dim F,\qquad q_F=d-k_F,\qquad H_F=\operatorname{lin}(F-F).
\]
The tangent cone at a point of $\relint F$, translated to the origin, is
denoted by $K_F$. It is a full-dimensional convex polyhedral cone with
lineality $H_F$. Its unit-ball volume is
\[
 \rho_F=|K_F\cap B_1(0)|>0.
\]
For the interior face, $K_A=\R^d$ and $\rho_A=\omega_d$, the volume of the
unit ball. Active facet inequalities have the form $a_i\cdot v\geq0$, with
common kernel $H_F$. Thus $K_F+H_F=K_F$, and
$K_F\cap H_F^\perp$ is full-dimensional in the $q_F$-dimensional normal space.
If $z,z'\in\aff F$, then $z+K_F=z'+K_F$. We may therefore use a fixed origin
in $\aff F$ in all normal-coordinate integrals.

\subsection{Face territories}
We recall the face territories of \cite[results~3.13--3.17]{PYH} in a form
needed below. For every fixed $S\geq1$ and all sufficiently small $r>0$, there
is a deterministic partition $\{T_F(r)\}$ of $A$ with the following properties:
\begin{enumerate}
\item If $x\in T_F(r)$, then $\dist(x,F)\leq D_F r$ for a fixed $D_F$.
Its tangential projection lies in a bounded set independent of $r$.
\item For $x\in T_F(r)$ and any $z\in\aff F$,
\begin{equation}\label{eq:localcone}
 A\cap B_{Sr}(x)=(z+K_F)\cap B_{Sr}(x).
\end{equation}
\item For each fixed $\eta>0$, the number of grid cubes of side $\eta r$
meeting $T_F(r)$ is $O(r^{-k_F})$.
\end{enumerate}
Here the territories and their constants depend on $S$.

For clarity, the construction uses the finite constant $K(A)$ in
\cite[Lemma~3.13]{PYH}. Choose $K_1=1$ and
$K_{j+1}>(2K(A)+1)K_j$. At scale $u=Sr$, the territory of a proper face $F$ is
\[
 F_{uK_{d-k_F}}\setminus
 \bigcup_{F'\text{ a proper face of }F}F'_{uK_{d-k_{F'}}},
\]
and the interior territory is the remainder. These territories partition
$A$ for small $u$. The published construction bounds their distance to $F$
and excludes every facet not containing $F$ from $B_u(x)$.
This proves \eqref{eq:localcone}: a segment from $x$ to a point of the cone
inside that ball could leave $A$ only through an inactive facet. The active
inequalities hold on the whole segment. Conversely, the active inequalities
hold throughout $A$.

A cube meeting $T_F(r)$ lies in a tube of thickness $O(r)$ about the bounded
$k_F$-dimensional face. The tube has volume $O(r^{q_F})$, whereas the cube has
volume $(\eta r)^d$, proving the cube count. Orthogonal normal coordinates and
Tonelli's theorem also give, for $\lambda>0$,
\begin{equation}\label{eq:depthintegral}
 \int_{T_F(r)}e^{-\lambda\dist(x,\aff F)}\,dx
 \leq c\lambda^{-q_F}.
\end{equation}
Indeed the tangential factor is bounded and the normal integral is bounded
by that over $\R^{q_F}$. Dilation reduces it to the finite integral of
$e^{-\|w\|}$; a zero-dimensional integral equals one.

\subsection{Three cone estimates}
\begin{lemma}\label{lem:translate}
Let $C\subset\R^d$ be a fixed full-dimensional closed convex cone and let
$\rho=|C\cap B_1(0)|$. For each fixed $R>0$ there is $c>0$ such that,
whenever $x,y\in C$ and $\|x-y\|\leq R$,
\[
 |(B_1(x)\cup B_1(y))\cap C|\geq\rho+c\min(\|x-y\|,1).
\]
\end{lemma}
\begin{proof}
The convex body $K=C\cap B_1(0)$ contains a ball of some radius $a>0$.
Also $x+K\subset C\cap B_1(x)$ and $y+K\subset C\cap B_1(y)$.
Take fibres parallel to $y-x$. Above a $(d-1)$-dimensional disk of radius
$a/2$ through the inscribed ball's centre, the fibres of $K$ are intervals of
length at least $a$. Translating an interval by $s=\|x-y\|$ increases the
length of its union with its translate by the smaller of $s$ and its length.
Integration gives
\[
 |(x+K)\cup(y+K)|\geq |K|+
 \omega_{d-1}(a/2)^{d-1}\min(s,a).
\]
Decreasing the constant proves the claim; $x=y$ is immediate. The direction
of translation need not lie in the cone.
\end{proof}

\begin{lemma}\label{lem:depth}
Let $C$ be a fixed full-dimensional convex polyhedral cone with lineality
$H$. For each fixed $D<\infty$ with $D\geq0$ there is $c_D>0$ such that
\[
 |B_1(x)\cap C|\geq |B_1(0)\cap C|+c_D\dist(x,H)
 \quad(x\in C,\ \dist(x,H)\leq D).
\]
\end{lemma}
\begin{proof}
The claim is immediate if $C=\R^d$ or $D=0$. Otherwise write an irredundant
representation $C=\{v:a_i\cdot v\geq0,\ 1\leq i\leq m\}$ with unit inward
facet normals. Set $w=\operatorname{proj}_{H^\perp}x$. Then
$w\in C\cap H^\perp$. On its compact set of unit vectors,
$\max_i a_i\cdot w$ has a positive minimum: if all products vanished, the
unit vector would lie in $H$. Hence some $i$ satisfies
$a_i\cdot w\geq c\|w\|$.

For each facet choose a positive-area closed disk $Q_i$ in its relative
interior and inside $B_{1/2}(0)$. All other facet inequalities have a positive
margin on $Q_i$; when there is only one facet this condition is empty.
For $v\in Q_i$ and $0<t<a_i\cdot w$, the point $v-t a_i$ is outside $C$,
but is in $C-w$ when $\|w\|$ is sufficiently small. In fact its $i$th
inequality after adding $w$ is nonnegative, and each other one is at least
its margin minus $\|w\|$. These points also lie in $B_1(0)$.
The resulting prism has volume $|Q_i|_{d-1}(a_i\cdot w)$.
Since $C\subset C-w$, this proves a linear gain for all sufficiently small
$\|w\|$, uniformly over the finitely many facets.

For larger $\|w\|\leq D$, the sets $C-tw$ increase as $t\geq0$ increases.
The gain at a fixed small displacement in the direction of $w$ supplies a
positive lower bound, which is at least a constant times $\|w\|/D$.
Finally $C+H=C$, so translating the ball centre by its tangential component
does not change the volume. This gives the asserted estimate.
\end{proof}

\begin{lemma}\label{lem:collar}
For a fixed full-dimensional closed convex cone $C$ and fixed
$0<\epsilon<R<\infty$, there is $\gamma>0$ such that every nonempty compact
$E\subset C$ with $\epsilon r\leq\diam E\leq Rr$ satisfies
\[
 |((E+B_r(0))\cap C)\setminus E|
 \geq (|C\cap B_1(0)|+\gamma)r^d.
\]
\end{lemma}
\begin{proof}
Scale to $r=1$ and put $K=C\cap B_1(0)$, $\rho=|K|$.
A diameter pair in $E$ and Lemma~\ref{lem:translate} give
$|E+K|\geq\rho+g$ for a fixed $g>0$. If $|E|\leq g/2$, subtracting $E$
leaves at least $\rho+g/2$. If $|E|>g/2$, the full Brunn--Minkowski inequality
for nonempty compact sets \cite[Theorem~4.1 and equation~(11)]{Gardner} gives
\[
 |E+K|-|E|\geq
 (|E|^{1/d}+\rho^{1/d})^d-|E|
 \geq\rho+d\rho^{(d-1)/d}(g/2)^{1/d}.
\]
Since $E+K\subset C\cap(E+B_1(0))$, these two cases give a common positive
surplus. No convexity or connectedness of $E$ is needed, and $E$ may have
zero volume.
\end{proof}

\section{A lower radius and three local estimates}
Define the face-dependent quantities
\begin{equation}\label{eq:critical}
 \beta=\max_F\frac{k_F}{d\rho_F},\quad
 \mathcal I=\left\{F:\frac{k_F}{d\rho_F}=\beta\right\},\quad
 b_F=\frac{(d-k_F)(d-1)}d,\quad
 a_* =\min_{F\in\mathcal I}\frac{b_F}{\rho_F}.
\end{equation}
The interior face implies $\beta>0$. Thus every critical face
$F\in\mathcal I$ has $k_F>0$. Fix $\varepsilon>0$ so small that
$\varepsilon\max_F\rho_F<(d-1)/2$, and put $a=a_*+\varepsilon$.
For all sufficiently large $n$, let
\begin{equation}\label{eq:lower}
 t_n=\beta\log n-a\log\log n,\qquad
 r_n=\left(\frac{t_n}{nf_0}\right)^{1/d}.
\end{equation}
Then $t_n>0$, $r_n\to0$, and $t_n/\log n\to\beta$.
The quantity
\begin{equation}\label{eq:Q}
 Q_F(n)=n^{k_F/d}t_n^{-b_F}e^{-\rho_Ft_n}
\end{equation}
will record the contribution of face $F$. For critical faces,
$Q_F(n)\asymp(\log n)^{a\rho_F-b_F}$ and
$a\rho_F-b_F\leq\varepsilon\rho_F$. For noncritical faces, $Q_F(n)$ is a
negative power of $n$ times a fixed power of $\log n$ and tends to zero.
For a critical face attaining the minimum in $a_*$,
$Q_F(n)\asymp(\log n)^{\varepsilon\rho_F}\to\infty$.

\subsection{Small components}
\begin{proposition}\label{prop:small}
There is a fixed $\delta\in(0,1)$ such that, with probability tending to one,
$G(\mathcal X_n,r_n)$ has no non-singleton component of diameter at most
$\delta r_n$.
\end{proposition}
\begin{proof}
First let $r$ be small with $t=nf_0r^d\asymp\log n$.
If $x$ is a vertex of such a component and $y$ a furthest vertex from $x$,
put $s=\|y-x\|\leq\delta r$. All component vertices lie in $B_s(x)$ and all
other sample points avoid
\[
 D(x,y)=\big[(B_r(x)\cup B_r(y))\setminus B_s(x)\big]\cap A.
\]
Use the territories with $S=3$. For $x\in T_F(r)$, the configuration agrees
with the face cone. Write $h_x=\dist(x,\aff F)$; its ratio to $r$ is bounded.
Scaling Lemmas~\ref{lem:translate} and \ref{lem:depth}, and taking half the
smaller gain, yields
\[
 |(B_r(x)\cup B_r(y))\cap A|
 \geq\rho_Fr^d+c r^{d-1}(h_x+s).
\]
The same $c>0$ works for all faces after taking a finite minimum.
Subtracting $B_s(x)$ loses at most $\omega_ds^d$. Choose $\delta$ sufficiently
small to absorb this loss in the term proportional to $s$. Thus
\begin{equation}\label{eq:smallvoid}
 |D(x,y)|\geq\rho_Fr^d+c' r^{d-1}(h_x+s).
\end{equation}

Count every ordered labelled pair as a possible witness. Conditional on
that pair, the other $n-2$ vertices avoid $D(x,y)$ with probability at most
$\exp[-(n-2)f_0|D(x,y)|]$. The expected number of witnesses is bounded by
\[
 n(n-1)f_0^2\sum_F\int_{T_F(r)}
 \int_{A\cap B_{\delta r}(x)}
 e^{-(n-2)f_0|D(x,y)|}\,dy\,dx.
\]
Apply \eqref{eq:smallvoid} on this domain before enlarging the integration
region of its exponential majorant to $\R^d$. The resulting inner integral
is $O((nr^{d-1})^{-d})$.
The depth integral \eqref{eq:depthintegral} contributes
$O((nr^{d-1})^{-q_F})$. Replacing $n-2$ by $n$ in the baseline exponential
costs a bounded factor tending to one. Substitution of $r^d=t/(nf_0)$ gives
the bound
\begin{equation}\label{eq:smallcount}
 c\sum_F n^{k_F/d}t^{-b_F-(d-1)}e^{-\rho_Ft}.
\end{equation}
At $t=t_n$, its critical summands are
$O((\log n)^{\varepsilon\rho_F-(d-1)})=o(1)$, and its noncritical summands
also vanish. Markov's inequality completes the proof. This count includes
every possible furthest pair, so it does not condition on an adaptively
chosen centre.
\end{proof}

\subsection{Medium separated subsets}
\begin{proposition}\label{prop:medium}
For each fixed $R>\delta$, with probability tending to one there is no
nonempty subset $U\subset\mathcal X_n$ with nonempty complement such that
\[
 \delta r_n\leq\diam U\leq Rr_n,
 \qquad \dist(U,\mathcal X_n\setminus U)>r_n.
\]
Connectedness of $U$ is not required.
\end{proposition}
\begin{proof}
Write $r=r_n$. Take a fixed grid of closed cubes of side $\eta r$, with
$\eta>0$ to be chosen. Let $B$ be the union of all cubes meeting $U$,
intersected with $A$. It is compact, contains $U$, and has diameter at most
$(R+2\sqrt d\eta)r$. Every point of $B$ lies within $\sqrt d\eta r$ of $U$.
Consequently the deterministic shell
\[
 H_B=\big[(B+B_{(1-\sqrt d\eta)r}(0))\cap A\big]\setminus B
\]
contains no sample point. A point outside $B$ in that shell would be within
$r$ of $U$, contradicting the assumed separation.

Choose $S=2(R+3)$ and initially $\eta\leq1/(2\sqrt d)$.
An anchor $x\in U$ belongs to one territory $T_F(r)$. The whole set $B$ and
its shell lie in $B_{Sr}(x)$ and hence agree with the local cone.
Set $s_0=1-\sqrt d\eta\in[1/2,1]$. Lemma~\ref{lem:collar}, at radius
$s_0r$ and with dimensionless diameter bounds $\delta$ and $2(R+2)$,
gives a surplus $\gamma_F>0$ independent of $\eta$ in this range.
Since there are finitely many faces, decrease $\eta$ so that for some
$\gamma'>0$ and every $F$,
\[
 (\rho_F+\gamma_F)s_0^d\geq\rho_F+\gamma'.
\]
It follows that $|H_B|\geq(\rho_F+\gamma')r^d$.

For a fixed anchor cube, only finitely many shapes for $B$ are possible,
with a bound independent of $n$: every cube in the shape lies within a
fixed multiple of $r$ of that cube. There are $O(r^{-k_F})$ anchor cubes
meeting $T_F(r)$. In this enumeration retain only shapes with the required
diameter bounds, containing the anchor cube after intersection with $A$,
and meeting that territory. The volume estimate applies to each such
deterministic shape by choosing a point of its intersection with $T_F(r)$.
Overlaps on grid boundaries add only a bounded number of possible cubes.
The binomial void bound and a union bound now give
\begin{equation}\label{eq:mediumvoid}
 \Pp(\text{a stated separated subset exists})
 \leq c\sum_F r^{-k_F}e^{-(\rho_F+\gamma')t_n}=o(1).
\end{equation}
Indeed $\beta\rho_F\geq k_F/d$, so every term is at most a fixed power of
$\log n$ times $n^{-\beta\gamma'}$. The estimate concerns all separated
subsets at one deterministic scale. It will therefore apply to components
at every larger radius.
\end{proof}

\subsection{An isolated vertex at the lower radius}
\begin{proposition}\label{prop:isolate}
The radius in \eqref{eq:lower} satisfies $\Pp(L_n>r_n)\to1$.
\end{proposition}
\begin{proof}
Choose a critical face $F$ with $b_F/\rho_F=a_*$. Its dimension is positive.
Let $P\subset\relint F$ be a compact patch of positive $k_F$-dimensional
volume, away from the relative boundary; for $F=A$ use an interior patch.
With $C_\perp=K_F\cap H_F^\perp$, define
\[
 S_n=\{z+w:z\in P,\ w\in C_\perp,\ \|w\|\leq r_n/t_n\}.
\]
If $q_F=0$, the normal factor has zero-dimensional volume one.
For large $n$, $S_n\subset A$ and radius-$3r_n$ neighbourhoods of its points
agree with the face cone. Its volume is
$v_n=c_P(r_n/t_n)^{q_F}$, with $c_P>0$.

Put $r=r_n$ and $t=t_n$. Cone inclusion gives the lower bound below.
For the upper bound, translating a ball by $w$ changes its intersection
volume by at most $c r^{d-1}\|w\|$: integrate interval fibres parallel to
$w$, whose changed length is at most a constant times $\|w\|$.
Uniformly for $x\in S_n$,
\[
 \rho_Fr^d\leq|B_r(x)\cap A|\leq\rho_Fr^d+c r^d/t.
\]
Let $Z_n$ count vertices in $S_n$ isolated at radius $r$. Given $X_i=x$,
its isolation probability is
$p(x)=(1-u(x))^{n-1}$, where $u(x)=f_0|B_r(x)\cap A|$.
We have $u=O(t/n)$ and $nu\in[\rho_Ft,\rho_Ft+O(1)]$.
The elementary bound
\[
 |\log(1-u)+u|\leq2u^2\quad(0\leq u\leq1/2)
\]
follows by integrating $v/(1-v)$ from $0$ to $u$.
Since $nu^2=O(t^2/n)=o(1)$, it implies
\begin{equation}\label{eq:isolatemean}
 \lambda_n=\E Z_n=nf_0\int_{S_n}p(x)\,dx
 \asymp n v_n e^{-\rho_Ft}
 \asymp Q_F(n)\asymp(\log n)^{\varepsilon\rho_F}\longrightarrow\infty.
\end{equation}

For the second factorial moment condition on two specified vertices
$x,y\in S_n$. At distance at most $r$ they cannot both be isolated.
For $r<\|x-y\|\leq2r$, Lemma~\ref{lem:translate} and the local cone
give a fixed $c_1>0$ such that their two ball neighbourhoods have union
volume at least $(\rho_F+c_1)r^d$. The contribution from these ordered pairs
is at most
\[
 c n^2v_n r^d e^{-(\rho_F+c_1)t}
 \leq c\lambda_n(nr^d)e^{-c_1t}=o(\lambda_n^2).
\]
Here $nr^d=O(\log n)$, whereas $e^{-c_1t}$ is a negative power of $n$
times a fixed logarithmic factor. The use of $n-2$ rather than $n$ in the
conditional void bound costs a factor tending to one.

At distance greater than $2r$ the balls are disjoint. The joint isolation
probability is $(1-u(x)-u(y))^{n-2}$. The logarithm bound gives, uniformly,
\[
 \frac{(1-u(x)-u(y))^{n-2}}{p(x)p(y)}=1+o(1),
\]
since $u(x),u(y)=O(\log n/n)$. Thus their contribution is at most
$(1+o(1))\lambda_n^2$. Combining the two ranges yields
$\operatorname{Var}Z_n=o(\lambda_n^2)$. Chebyshev's inequality proves
$\Pp(Z_n=0)\to0$. An isolated vertex at the closed radius $r$ has
nearest-neighbour distance strictly greater than $r$, as required.
\end{proof}

\section{Large components at all larger radii}
We first record a deterministic separation argument. We use two geometric
facts from \cite[Lemma~3.14 and Proposition~3.15]{PYH}: $A$ is unicoherent,
and there are constants $\delta_1>0$ and $K_0>1$ such that, for $x\in A$
and $0<b<\delta_1$, $A\setminus B_b(x)$ has at most two components; if
there are two, one has diameter at most $K_0b$. Unicoherence means that the
intersection of two closed connected sets whose union is $A$ is connected.
The following argument expands the deterministic part of
\cite[proof of Lemma~3.12, equation~(3.18)]{PYH}; no probabilistic conclusion
of that lemma is used.

\begin{lemma}\label{lem:separator}
There is $c_A>0$ such that, whenever nonempty finite sets $U,W\subset A$
satisfy $\dist(U,W)>s>0$ and both $G(U,s)$ and $G(W,s)$ are connected,
there is a nonempty compact connected set $S\subset A$ separating every
continuous path from $U$ to $W$ and satisfying
\[
 \dist(z,U)=s/2\quad(z\in S),\qquad
 \diam S\geq c_A\min(\delta_1,\diam U,\diam W).
\]
\end{lemma}
\begin{proof}
Put $K=U_{s/2}$. It is compact and connected: its finitely many convex ball
sections connect along graph edges, whose midpoints lie in $A$.
Similarly $W_{s/2}$ is connected and disjoint from $K$.
The relative open set $A\setminus K$ is locally path connected, because a
sufficiently small relative ball about any of its points is convex and stays
in the set. Its components are consequently open and path connected.
Let $V$ be the component containing $W_{s/2}$, and set $U'=A\setminus V$.

Every component $Q$ of $A\setminus K$ has boundary in $K$ and its closure
meets $K$. The first assertion follows from its relative closedness in
$A\setminus K$ and the openness of the other components. Otherwise $Q$
would be a nonempty proper clopen subset of connected $A$.
Hence $U'$ is the union of $K$ and the closures of all components other
than $V$, each attached to $K$, and is connected. Both $U'$ and
$\overline V$ are closed and connected, and their union is $A$.
Unicoherence makes $S=\overline V\cap U'$ compact and connected.
It is nonempty and every path from $U$ to $W$ crosses it. Every point of
$S$ lies in $K$ and has distance exactly $s/2$ from $U$: a point of smaller
distance has a relative neighbourhood contained in $K$, disjoint from $V$.

Set $m=\min(\delta_1,\diam U,\diam W)$ and
$c_A=1/[6(K_0+1)]$. If $m=0$ the diameter bound is immediate.
Otherwise, if it failed, choose a positive $b$ with
$\diam S<b<c_A m$ and $z_0\in S$. Then $S\subset B_b(z_0)$ and both
$U$ and $W$ have diameter greater than $6(K_0+1)b$.
Choose $u\in U$ and $w\in W$ outside that ball. If its complement in $A$
had one component, a path there would join $u$ to $w$ without meeting $S$.
Thus it has two components, with $u,w$ in different ones.
One, say that containing $u$, has diameter at most $K_0b$.
A path in convex $A$ from $u$ to $w$ first exits this component at a point
of its closure in $B_b(z_0)$, so $\dist(u,B_b(z_0))\leq K_0b$.

There is $u'\in U$ with $\|u-u'\|>(2K_0+2)b$; otherwise
$\diam U\leq(4K_0+4)b$, a contradiction. This point is outside both the
small component and the ball: points of the ball are at distance at most
$(K_0+2)b$ from $u$. It therefore lies in the component containing $w$,
and a path there again contradicts separation by $S$. This proves the
diameter bound, including a possible zero-diameter separator by the choice
of a positive enclosing radius $b$.
\end{proof}

We also need uniform local volume: there is $c_0>0$ such that
\begin{equation}\label{eq:localmass}
 |A\cap B_h(x)|\geq c_0h^d\qquad(x\in A, h\text{ sufficiently small}).
\end{equation}
Here is a direct proof. Fix $B_e(c)\subset A$ with $e>0$ and
$D\geq\sup_{x\in A}\|c-x\|$. For $\alpha=h/(D+e)<1$, convexity places
$(1-\alpha)x+\alpha B_e(c)$ inside both $A$ and $B_h(x)$.
Its volume is $\omega_d[eh/(D+e)]^d$.

\begin{proposition}\label{prop:large}
For some fixed $R<\infty$, with probability tending to one no graph
$G(\mathcal X_n,s)$ with $s\geq r_n$ has two distinct components each of
diameter at least $Rr_n$.
\end{proposition}
\begin{proof}
Write $r=r_n$. Fix $\eta\in(0,1/20)$ and $h=\eta r$, and take a deterministic
maximal $h$-separated net $\mathcal L\subset A$.
Its $h$-balls cover $A$; disjoint ambient $h/3$-balls show that its cardinality
$N$ is $O(h^{-d})$. The graph on $\mathcal L$ joining sites at distance at
most $2h$ has degree bounded in terms of $d$ alone. For a fixed constant
$D_0$, the number of its connected $q$-site subsets is at most $ND_0^q$.
Indeed each can be traversed by a depth-first walk of length $2(q-1)$,
and counting all such walks overcounts the subsets.

Suppose two components $U,W$ of $G(\mathcal X_n,s)$ have diameter at least
$Rr$ for some $s\geq r$. For large $n$, $Rr<\delta_1$.
Lemma~\ref{lem:separator} supplies $S$ with $\diam S\geq c_ARr$.
Let
\[
 D_S=\{z\in\mathcal L:\dist(z,S)\leq h\}.
\]
Its radius-$h$ ball sections cover $S$, are connected, and each meets $S$.
Their union is connected, so $D_S$ is connected in the net graph.
A connected union of $q$ balls of radius $h$ has diameter at most $2hq$;
therefore $q=|D_S|\geq c_AR/(2\eta)$.

The set
\[
 V_{D_S}=\bigcup_{z\in D_S}(B_{h/3}(z)\cap A)
\]
contains no sample point. For a point $v$ there, the Lipschitz property of
distance to $U$ and the separator identity give
\[
 s/2-4h/3\leq\dist(v,U)\leq s/2+4h/3.
\]
Since $s\geq r$ and $\eta<1/20$, these bounds lie strictly between $0$ and
$s$. Such a sample point would neither belong to $U$ nor be allowed outside
its graph component.

For every deterministic $q$-site set $D\subset\mathcal L$, the ambient
$h/3$-balls are disjoint and \eqref{eq:localmass} yields
$|V_D|\geq c_2qh^d$. Its void probability is at most
$\exp(-nf_0c_2qh^d)$. A union bound over the one fixed net gives
\begin{equation}\label{eq:largevoid}
 \Pp(\text{two stated large components exist at some }s\geq r_n)
 \leq N\sum_{q\geq c_AR/(2\eta)}D_0^q e^{-nf_0c_2qh^d}.
\end{equation}
Now $nf_0h^d=\eta^dt_n\geq(\eta^d\beta/2)\log n$ for large $n$, and
$N=O(n/\log n)$. Choose a fixed $R$ large enough that the lower bound on
$q$ makes this geometric series, including $N$, $O(n^{-2})$.
This proves the result simultaneously for all $s\geq r_n$.
There is no union bound over a continuum of radii: each separator produces
an empty subset of the same deterministic net. The event is measurable
because a finite sample graph changes only at finitely many pair distances.
\end{proof}

\section{Proof of the hitting-time equality}
Take $\delta$ from Proposition~\ref{prop:small} and enlarge $R$ from
Proposition~\ref{prop:large} if necessary so that $R>\delta$.
Apply Proposition~\ref{prop:medium} with this fixed $R$.
On the intersection of these three events, fix any $s\geq r_n$ for which
$G(\mathcal X_n,s)$ is disconnected.

A non-singleton component $U$ cannot have diameter less than $\delta r_n$.
Its vertices would already form a complete graph at $r_n$, since $\delta<1$,
and would have no neighbours outside $U$ at $r_n$. This would contradict
Proposition~\ref{prop:small}. Nor can its diameter lie between $\delta r_n$
and $Rr_n$: its distance from its nonempty complement is strictly greater
than $s\geq r_n$, contrary to Proposition~\ref{prop:medium}.
Thus every non-singleton component has diameter greater than $Rr_n$.
Proposition~\ref{prop:large} allows at most one such component. A disconnected
graph must therefore have an isolated vertex.

By Proposition~\ref{prop:isolate}, also $L_n>r_n$ with probability tending to
one. At the closed radius $s=L_n$, every vertex has a neighbour by
\eqref{eq:thresholds}. The uniform conclusion just proved forces
$G(\mathcal X_n,L_n)$ to be connected. Hence $M_n\leq L_n$, and the
deterministic reverse inequality gives $L_n=M_n$.
The four propositions each fail with probability tending to zero, so their
finite intersection proves Theorem~\ref{thm:main}.

The conclusion is exact equality with high probability for every fixed
$d\geq3$ and every fixed full-dimensional compact convex polytope.
The constants in the auxiliary radius and cone estimates depend on that
polytope; no uniform claim is made when its angles or shape vary.

\subsection*{AI use statement}
AI assistance through OpenAI Codex was used to develop and write the
argument, check mathematical and bibliographic dependencies, and prepare
this manuscript. The resulting solution claim requires independent
mathematical review.

\begin{thebibliography}{9}
\bibitem{Gardner}
R. J. Gardner, \emph{The Brunn--Minkowski inequality},
Bulletin of the American Mathematical Society (New Series)
\textbf{39} (2002), no.~3, 355--405.
\href{https://doi.org/10.1090/S0273-0979-02-00941-2}{doi:10.1090/S0273-0979-02-00941-2}.

\bibitem{PYH}
M. D. Penrose, X. Yang and F. Higgs,
\emph{Largest nearest-neighbour link and connectivity threshold in a
polytopal random sample}, Journal of Applied and Computational Topology
\textbf{8} (2024), 1723--1750.
\href{https://doi.org/10.1007/s41468-023-00154-5}{doi:10.1007/s41468-023-00154-5}.

\bibitem{PY}
M. D. Penrose and X. Yang,
\emph{Fluctuations of the connectivity threshold and largest
nearest-neighbour link}, arXiv:2406.00647v3 (2024).
\href{https://arxiv.org/abs/2406.00647v3}{https://arxiv.org/abs/2406.00647v3}.
\end{thebibliography}
\end{document}
